\chapter{Studi Literatur}
\label{chapter:studi-literatur}

Bab Studi Literatur digunakan untuk mendeskripsikan kajian literatur yang terkait dengan persoalan tugas akhir. Tujuan studi literatur adalah:

\begin{enumerate}
	\item menunjukkan kepada pembaca adanya gap seperti pada rumusan masalah yang memang belum terselesaikan,
	\item memberikan pemahaman secukupnya kepada pembaca tentang teori atau pekerjaan terkait yang terkait langsung dengan penyelesaian persoalan, serta
	\item menyampaikan informasi apa saja yang sudah ditulis/dilaporkan oleh pihak lain (peneliti/Tugas Akhir/Tesis) tentang hasil penelitian/pekerjaan mereka yang sama atau mirip kaitannya dengan persoalan tugas akhir.
\end{enumerate}

Kita juga dapat memisah beberapa bagian latex untuk \textit{readability}. Kita dapat memasukan file latex lainnya dengan menggunakan fitur input. Berikut merupakan contoh ketika memasukan bagian contoh-subbab ke chapter-2

\section{Contoh Subbab}
Perujukan literatur dapat dilakukan dengan menambahkan entri baru di berkas.

Sekarang mau ke bab berapa yaaaa.... hmm... ke bab \ref{chapter:studi-literatur} ahhhhh.

\blindtext

\subsection{Contoh Subsubbab}
\label{subsec:contoh-subsec}

\blindtext

\begin{figure}[ht]
	\centering
	\includegraphics[width=0.5\textwidth]{resources/cover-ganesha.jpg}
	\caption{Contoh gambar}
\end{figure}

\subsubsection{Subsubsubbab}
\blindtext

\subsubsubsection{sub sub sub sub bab}
\blindtext

\begin{table}[h]
	\caption{Tabel random}
	\vspace{0.25cm}
	\begin{center}
		\begin{tabular}{|c|c|c|c|}
			\hline
			Title1 & Title2 & Title3 & Title4  \tabularnewline
			\hline
			1647   & 1.97   & 0.68   & 1.90 \tabularnewline
			2301   & 2.92   & 1.06   & 2.75 \tabularnewline
			2969   & 3.23   & 1.16   & 3.78 \tabularnewline
			3791   & 4.39   & 1.40   & 4.14 \tabularnewline
			4625   & 6.72   & 1.87   & 5.59 \tabularnewline
			\hline
		\end{tabular}
	\end{center}
\end{table}  % Now fixed - citations added to bibliography

\section{Menyisipkan Persamaan}

Beberapa contoh menyisipkan persamaan.

\subsection{Contoh Bikin Equation}
\textbf{text tebal} dan ini \emph{miring}, bikin persamaan di baris yang sama, tinggal pake dolar2 $\Psi(\vec{r}_1,...,\vec{r}_N)$, sehingga persamaan Schr\"{o}dinger, terus, persamaan yang dinomeri kayak gini
%ini contoh bikin persamaan, ..... :D
\begin{equation}
	\left[ \sum_{i}^{N}-\frac{\hbar^2}{2m}\nabla_i^2 + \sum_{i}^{N}V(\vec{r}_i)+ \sum_{i<j}^{N}(\vec{r}_i,\vec{r}_j)\right]\Psi = E\Psi
\end{equation}

untuk $N$-elektron, dengan $\hat{H}$=Hamiltonian, $E$=Energi total, $\hat{T}$=Energi kinetik, $\hat{V}$=Energi potensial, dan $\hat{U}$=Interaksi ektron-elektron.

\subsection{Bikin Matrix}
Lalalallala.... bikin matrix sekarang, yang ini dikecilin, pake smaller
	{\smaller
		\begin{equation}
			\Psi({\bf r}_1, {\bf r}_2, \cdots {\bf r}_N) = \frac{1}{\sqrt{N!}}\left| \begin{array}{llcl}
				\phi_1({\bf r}_1)     & \phi_2({\bf r}_1)     & \cdots                & \phi_N({\bf r}_1)     \\
				\phi_1({\bf r}_2)     & \phi_2({\bf r}_2)     & \cdots                & \phi_N({\bf r}_2)     \\
				\phi_1({\bf r}_3)     & \phi_2({\bf r}_3)     & \cdots                & \phi_N({\bf r}_3)     \\
				\multicolumn{1}{c}{.} & \multicolumn{1}{c}{.} & \multicolumn{1}{c}{.} & \multicolumn{1}{c}{.} \\
				\multicolumn{1}{c}{.} & \multicolumn{1}{c}{.} & \multicolumn{1}{c}{.} & \multicolumn{1}{c}{.} \\
				\multicolumn{1}{c}{.} & \multicolumn{1}{c}{.} & \multicolumn{1}{c}{.} & \multicolumn{1}{c}{.} \\
				\phi_1({\bf r}_N)     & \phi_2({\bf r}_N)     & \cdots                & \phi_N({\bf r}_N)     \\
			\end{array} \right|
		\end{equation}
	}